\documentclass{article}
\usepackage[margin=1in]{geometry}
\usepackage{amsmath}
\usepackage{cite}

%--- for references ---------
\usepackage[
    backend=biber,
    style=numeric,
    sorting=none
    ]{biblatex}

\addbibresource{bibliography.bib}
\AtBeginBibliography{\vspace*{10pt}}
%----------------------------

% Parskip and parindent
\setlength{\parindent}{0pt} % Begin of paragraph indentation
\setlength{\parskip}{1em} % Paragraph spacing


\begin{document}

\begin{center}
    \LARGE Self-Supervised Learning for Video Geometry\\[0.7em]  % Adjust spacing here
    \large Vanja Kovinic
    \vspace{1cm}
\end{center}

\section{Introduction}

Recent advances in 3D vision have produced feed-forward models that reconstruct geometry without optimization. VGGT \cite{vggt} predicts camera poses, depth, and structure from multi-view images in under one second; Geo4D \cite{geo4d} extends to dynamic scenes using video diffusion priors. Both rely on labeled or synthetic supervision, limiting scalability to billions of hours of unlabeled video online.

Videos offer an extraordinarily rich yet underutilized signal for geometric learning. Unlike static images, videos provide: (1) natural multi-view geometry through camera motion, (2) temporal consistency constraints that implicitly encode 3D structure, (3) motion cues that reveal depth ordering and object boundaries, and (4) virtually unlimited scale - billions of hours of video exist without geometric labels. Self-supervised learning transformed language \cite{bert} and vision \cite{mae,dino} by exploiting unlabeled data. Can it do the same for video geometry?

Self-supervised monocular depth exists \cite{godard2019}, but methods focus on per-frame 
predictions. Designing objectives that teach how 3D 
structure evolves over time remains unexplored. Which temporal self-supervised objectives 
best capture 4D structure? How data-efficient is video pretraining? What geometric 
understanding emerges from unlabeled videos?

\section{Background}

\subsection{Feed-Forward Multi-View Reconstruction}

Multi-view reconstruction traditionally requires optimization-based structure-from-motion 
\cite{sfm}. Recent work replaces this with feed-forward prediction. 
DUSt3R \cite{DUSt3R} pioneered dense point regression from image pairs but requires 
global alignment for multiple views. VGGT \cite{vggt} extends this with camera tokens 
and alternating attention, processing 1-200+ images jointly. MonST3R \cite{MonST3R} 
performs pairwise tracking for dynamic scenes; Geo4D \cite{geo4d} adapts video diffusion 
models. All require expensive supervision: DUSt3R/VGGT use structure-from-motion labels, 
Geo4D uses synthetic data rendered from 3D assets, limiting real-world scalability.

\subsection{Self-Supervised Learning}

Self-supervised learning extracts training signals from data structure. Masked 
autoencoders \cite{mae} reconstruct masked patches; DINO \cite{dino,dinov2} matches 
predictions between augmented views, producing semantic features without labels. 
VideoMAE \cite{videomae} extends masking to spatiotemporal tubes. For geometry, 
self-supervised depth \cite{godard2019,watson2021} uses photometric consistency between 
consecutive frames. However, methods predict per-frame depth independently, treating 
video as image sequences rather than modeling temporal geometric reasoning. Whether 
self-supervised objectives can teach spatiotemporal structure (tracking evolving 3D 
geometry and maintaining multi-view consistency) remains unexplored, yet is fundamental 
to scaling video geometric understanding.

\newpage
\printbibliography[title={References}]
\end{document}